% Recuperación expositiva de INTERRELACIONES_ESTRUCTURALES_VACIO.md, §§6--8.
% Insertar después de la inversión cúbica y antes de las coordenadas logarítmicas.
\subsection{Composición transportada de las respuestas constitutivas}
\label{amp-comp-sec}

El transporte angular y las incidencias $90$ y $120$ proceden de la
misma estructura APP--TRIT--TPK: las hojas conservan las evaluaciones y
sus acarreos; el régimen trítico conserva la orientación; el calendario
TPK transporta fase y memoria hasta los canales angulares. La respuesta
constitutiva aplica a esos canales la función $f$ de
\eqref{iii:eq:monotone}. Su inversión permite expresar, en las coordenadas
de respuesta, la composición de transportes del mismo marco. Se conserva
así el orden constructivo: transporte, composición y evaluación espectral.
La operación obtenida actúa sobre esta representación del estado enriquecido;
su inversa recupera los canales, mientras la historia completa de memoria
permanece en la estructura de procedencia.

Sean $s=q^{30}$ y $\psi=f^{-1}$. El exponente $30$ es el máximo común
divisor de las incidencias $90$ y $120$. La inversión cúbica determina
$\psi:(3/4,1)\longrightarrow(0,1)$. Adjuntando la respuesta de frontera
$f(1)=3/4$, extendemos la biyección a
\[
 \psi:D\longrightarrow(0,1],\qquad D=[3/4,1).
\]
En esa frontera se utiliza la función racional $f$; el cociente
$(1-q^{90})/(1-q^{120})$ tiene allí una singularidad removible.

\begin{proposition}[Monoide de respuestas y coordenada aditiva]
\label{amp-comp-monoid}
La operación
\begin{equation}
 r\odot t=f\bigl(\psi(r)\psi(t)\bigr),\qquad r,t\in D,
 \label{amp-comp-law}
\end{equation}
hace de $D$ un monoide conmutativo y cancelativo con elemento neutro $3/4$.
La aplicación
\begin{equation}
 \ell(r)=-\frac1{30}\log\psi(r)
 \label{amp-comp-character}
\end{equation}
es un isomorfismo de ese monoide con $(\mathbb R_{\geq0},+)$.
El intervalo $(3/4,1)$ es cerrado bajo $\odot$; el neutro pertenece a
la extensión de frontera. El único elemento invertible dentro de $D$
es el neutro.
\end{proposition}
\begin{proof}
El producto de dos elementos de $(0,1]$ pertenece al mismo intervalo y
\[
 \psi(r\odot t)=\psi(r)\psi(t).
\]
Para tres respuestas, ambos órdenes de asociación tienen imagen
$\psi(r)\psi(t)\psi(u)$; la inyectividad de $\psi$ prueba la
asociatividad. La conmutatividad procede del producto escalar. Si
$r\odot t=r\odot u$, la positividad de $\psi(r)$ permite cancelar y
obtener $\psi(t)=\psi(u)$, luego $t=u$. El parámetro $1$ corresponde a
$3/4$ y proporciona la identidad. El logaritmo transforma el producto en
suma, y su rango en $(0,1]$ prueba que $\ell$ es una biyección sobre
$\mathbb R_{\geq0}$. En particular, para la respuesta de un canal $q$,
$\ell(r)=-\log q$. Si ambos parámetros son estrictamente menores que
$1$, su producto también lo es. Finalmente, dos números de $(0,1]$
tienen producto $1$ únicamente si ambos son $1$.
\end{proof}

\begin{proposition}[Cancelación admisible y prolongación inversa]
\label{amp-comp-inverse}
Para $r,u\in D$, la ecuación $r\odot t=u$ tiene solución $t\in D$
exactamente cuando $u\geq r$. Esa solución es única y viene dada por
\begin{equation}
 t=f\!\left(\frac{\psi(u)}{\psi(r)}\right).
 \label{amp-comp-cancellation}
\end{equation}
La misma función $f$ es una biyección de $(0,\infty)$ sobre $(0,1)$.
Con esta extensión de $\psi$, la ley~\eqref{amp-comp-law} transporta el
grupo multiplicativo de los parámetros positivos al intervalo $(0,1)$.
La inversión de grupo satisface
\begin{equation}
 r^{\ominus}=\psi(r)r,\qquad
 r\odot r^{\ominus}=\frac34.
 \label{amp-comp-group-inverse}
\end{equation}
\end{proposition}
\begin{proof}
La ecuación equivale a
$\psi(t)=\psi(u)/\psi(r)$. Este cociente pertenece a $(0,1]$
exactamente cuando $\psi(u)\leq\psi(r)$, que por monotonía decreciente
equivale a $u\geq r$. La biyección prueba existencia y unicidad.
La derivada de \eqref{iii:eq:monotone} es negativa para todo $s>0$;
los límites de $f$ en $0$ y en infinito son $1$ y $0$, respectivamente.
Esto establece la biyección extendida. Multiplicando numerador y
denominador por $s^3$ se obtiene
\[
 f(s^{-1})=\frac{s^3+s^2+s}{s^3+s^2+s+1}=s f(s).
\]
La inversa multiplicativa de $s=\psi(r)$ tiene, por tanto, respuesta
$\psi(r)r$. Su composición con $r$ corresponde al parámetro $1$.
\end{proof}

La prolongación a $s>1$, o equivalentemente a $q>1$, es la extensión
algebraica del lector. El sector físico contractivo del capítulo conserva
su dominio declarado. La coordenada $\ell$ se prolonga a un isomorfismo
con $(\mathbb R,+)$ y cambia de signo bajo la inversión de grupo.

\subsubsection{Realización operatoria en un marco de hojas común}

\begin{proposition}[Composición con proyectores marcados]
\label{amp-comp-operator}
Sean $\mathcal R=\mathcal R^*$, $\mathcal R^2=I$ y
$P_\pm=(I\pm\mathcal R)/2$. Para $j=a,b$, sean
\[
 T_j=q_{j,+}P_++q_{j,-}P_-,\qquad
 \mathcal V_j=f(T_j^{30}),\qquad 0<q_{j,\pm}<1.
\]
La composición $T_{a\circ b}=T_aT_b$ satisface
\begin{equation}
 \mathcal V_{a\circ b}
 =f\bigl(\psi(\mathcal V_a)\psi(\mathcal V_b)\bigr).
 \label{amp-comp-operator-law}
\end{equation}
Sus dos respuestas son $r_{a,\pm}\odot r_{b,\pm}$. Si
$T_j=\exp[-(x_jI+y_j\mathcal R)]$ con $x_j>|y_j|$, entonces
\begin{equation}
 T_aT_b=
 \exp[-((x_a+x_b)I+(y_a+y_b)\mathcal R)],
 \label{amp-comp-angular-addition}
\end{equation}
y $x_a+x_b>|y_a+y_b|$.
\end{proposition}
\begin{proof}
La ortogonalidad $P_+P_-=0$ da
\[
 T_aT_b=q_{a,+}q_{b,+}P_++q_{a,-}q_{b,-}P_-.
\]
En particular, $T_a$ y $T_b$ conmutan. El cálculo funcional en estos
proyectores proporciona
$\psi(\mathcal V_j)=T_j^{30}$ y
$(T_aT_b)^{30}=T_a^{30}T_b^{30}$; aplicar $f$ prueba
\eqref{amp-comp-operator-law}. Para toda función $g$ definida en los
dos autovalores se tiene, más explícitamente,
\[
 P_\pm g(T_j)=g(q_{j,\pm})P_\pm.
\]
Como $q_{j,\pm}=\exp[-(x_j\pm y_j)]$, sus productos prueban
\eqref{amp-comp-angular-addition}. La desigualdad triangular concluye
$x_a+x_b>|y_a|+|y_b|\geq|y_a+y_b|$.
\end{proof}

Las coordenadas recuperadas pueden escribirse directamente como
\[
 x=\frac{\ell(r_+)+\ell(r_-)}2,\qquad
 y=\frac{\ell(r_+)-\ell(r_-)}2.
\]
La composición suma estas parejas. El intercambio simultáneo de hojas
permuta los dos canales de cada transporte, actúa como
$(x,y)\mapsto(x,-y)$ y es un automorfismo de la ley componente a
componente. La inversión de ambos transportes en la extensión algebraica
actúa como $(x,y)\mapsto(-x,-y)$. Las dos involuciones conmutan y
conservan funciones distintas: una intercambia las marcas de orientación;
la otra invierte el transporte.

Para representar el intercambio mediante conjugación se utiliza un
operador $L$ que satisfaga $L\mathcal R L^{-1}=-\mathcal R$. En la
representación
\[
 \mathcal R=\begin{pmatrix}0&1\\1&0\end{pmatrix},\qquad
 L=\begin{pmatrix}1&0\\0&-1\end{pmatrix},
\]
la multiplicación directa da $LP_\pm L^{-1}=P_\mp$ y, por tanto,
$LT_jL^{-1}=q_{j,-}P_++q_{j,+}P_-$. La conjugación por
$\mathcal R$ conserva sus propios proyectores. Una traza que reúne los
canales tampoco reemplaza a los proyectores en esta identidad funcional:
la información de sus etiquetas interviene en la composición.

La hipótesis de marco común es sustantiva. Por ejemplo, los operadores
positivos definidos
\[
 A_0=\begin{pmatrix}1/2&0\\0&1/3\end{pmatrix},\qquad
 B_0=\begin{pmatrix}1/2&1/6\\1/6&1/2\end{pmatrix}
\]
tienen autovalores positivos, pero
\[
 A_0B_0=\begin{pmatrix}1/4&1/12\\1/18&1/6\end{pmatrix}
\]
carece de autoadjunción. La proposición especifica la composición de
transportes en los proyectores comunes; su dominio no comprende una mezcla
arbitraria de medios materiales. La realización de una concatenación de
rutas HMT conserva además sus datos de memoria: el cierre algebraico del
cono de parámetros describe la representación y no selecciona por sí solo
nuevas rutas primarias.

\subsubsection{Producto de canales y composición de transportes}

La velocidad normalizada de un estado sigue determinada por
\[
 \widehat c=(r_+r_-)^{-1}.
\]
Aquí el producto ordinario reúne las dos respuestas de ese mismo estado.
La ley $\odot$ evalúa otra operación: la composición de dos transportes
en cada hoja. El siguiente cálculo exacto distingue sus argumentos:
\begin{equation}
 f(1/2)=\frac{14}{15},\qquad
 \frac{14}{15}\odot\frac{14}{15}=f(1/4)=\frac{84}{85}
 \ne\frac{196}{225}=\left(\frac{14}{15}\right)^2.
 \label{amp-comp-rational-example}
\end{equation}
Las fracciones de este ejemplo son parámetros algebraicos de comprobación.
La expresión de $\widehat c$ y las identidades constitutivas anteriores
permanecen inalteradas. El contenido añadido es una ley explícita para
transportar la composición y recuperar su coordenada angular aditiva.
